% Do not forget to include Introduction
%---------------------------------------------------------------
\chapter{Introduction}
% uncomment the following line to create an unnumbered chapter
%\chapter*{Introduction}\addcontentsline{toc}{chapter}{Introduction}\markboth{Introduction}{Introduction}
%---------------------------------------------------------------
\setcounter{page}{1}

Variable stars play a fundamental role in advancing our understanding of the universe, particularly in cosmology and stellar physics.
Their study has contributed to numerous major scientific discoveries, such as the calculation of the rate of at which the universe expands. \parencite{Hubble1929}
In addition, the analysis of variable stars can support the detection of extraterrestrial planets which is one of the main areas of research at the Astronomical Institute of Czech Academy of Sciences.

Recent technological progress has led to substantial improvements in photometric instrumentation, enabling the development of large-scale astronomical surveys such as OGLE (Optical Gravitational Lensing Experiment --~\cite{Udalski2008}), the Kepler planet-detection mission \parencite{Borucki2010}, the Zwicky Transient Facility survey \parencite{Bellm2019}, and TESS (Transiting Exoplanet Survey Satellite --~\cite{Ricker2014}).
These surveys generate enormous volumes of high-quality astronomical data, including light curves.
However, the scale of these datasets makes manual classification of variable stars both time-consuming and impractical, creating a strong need for automated methods.

Machine learning approaches have proven to be an effective solution for automated classification, enabling researchers to process large amounts of data quickly and accurately.
Moreover, deep learning methods, such as CNNs, have often outperformed more traditional machine learning techniques, such as Random Forests and K-means clustering.

In recent years, transformer-based models have emerged as a promising alternative to earlier deep learning approaches in many domains, including time-series classification, where they often achieve very good performance.
Fine-tuning pretrained transformer models, such as LLMs, TSFMs and VTs, has also shown potential in a variety of tasks by leveraging the powerful pattern-recognition capabilities acquired during pretraining.

In this thesis, we investigate the viability of transformer-based models for light curve classification.
Although previous studies have already explored this topic, our goal is to provide a more comprehensive analysis by considering both transformers trained from scratch and fine-tuned large pretrained models.



